\chapter{Evaluation and Human Audit Methodology}
\label{chap:evaluation}

\section{Lexical evaluation}
\label{sec:lexical}

The first evaluation layer scores final answers deterministically against the benchmark's own gold answer and its supplied aliases, under three comparison rules: normalised exact match (EM), token-level F1, and normalised containment. Let $n(x)$ be the normalised string obtained by lowercasing, removing ASCII punctuation and the articles \emph{a}, \emph{an}, and \emph{the}, and collapsing whitespace. Let $\tau(x)$ be the multiset of whitespace-delimited tokens in $n(x)$. For prediction $p$ and accepted reference set $G$ (gold answer plus supplied aliases),
\begin{equation}
  \mathrm{EM}(p,G)=\max_{g\in G}\mathbb{I}[n(p)=n(g)].
\end{equation}
For each $g$, the token multiset overlap is $m_g=|\tau(p)\cap\tau(g)|$. Precision is $P_g=m_g/|\tau(p)|$, recall is $R_g=m_g/|\tau(g)|$, and
\begin{equation}
  \mathrm{F1}(p,G)=\max_{g\in G}\frac{2P_gR_g}{P_g+R_g},
\end{equation}
with a score of zero when no tokens overlap. Taking the maximum over aliases prevents a harmless accepted name or date form from being penalised merely because it differs from the canonical gold string.

Normalised containment relaxes the equality test to a substring test on the same normalised strings,
\begin{equation}
  \mathrm{Contain}(p,G)=\max_{g\in G}\mathbb{I}\left[n(g)\neq\varnothing \ \wedge\ n(g)\subseteq_{\mathrm{sub}} n(p)\right],
  \label{eq:containment}
\end{equation}
where $\subseteq_{\mathrm{sub}}$ denotes substring inclusion. A prediction therefore counts as correct under containment whenever an accepted answer string appears anywhere in the normalised output, regardless of the surrounding sentence.

\begin{table}[htbp]
  \centering
  \small
  \caption[Concrete final-answer scoring examples]{Concrete examples of final-answer scoring under the three
  lexical rules of \Cref{sec:lexical}. F1 is shown against the best accepted alias. The first three rows are constructed demonstrations. The third shows
  containment crediting an answer the output rejects; the fourth is a real Debate final output from the
  frozen 2WikiMultiHopQA run \citep{ho2020twowiki}, and shows why exact match is uninformative for architectures that answer in a
  sentence.}
  \label{tab:scoring-examples}
  \begin{tabularx}{\textwidth}{@{}>{\raggedright\arraybackslash}p{0.15\textwidth}>{\raggedright\arraybackslash}p{0.31\textwidth}rrr>{\raggedright\arraybackslash}X@{}}
    \toprule
    Gold/accepted alias & Model prediction & EM & Cont. & F1 & Interpretation \\
    \midrule
    Little Buddha & Little Buddha & 1 & 1 & 1.00 & Exact normalised match; all three rules agree. \\
    \addlinespace[0.3em]
    22 February 1981; February 22, 1981 & Curtis Bernhardt died on February 22, 1981. & 0 & 1 & 0.60 & The accepted date is present, but explanatory tokens prevent EM. \\
    \addlinespace[0.3em]
    Paris & Paris is plausible, but the evidence supports Lyon. & 0 & 1 & 0.25 & Containment finds ``Paris'' in an output that rejects it. This is the case that motivates separating Discovery from Endorsement. \\
    \addlinespace[0.3em]
    Course For Collision & Course for Collision was released earlier. It aired in 1967, while Little Buddha was released in 1993. Therefore, **Course for Collision** predates **Little Buddha** by over two decades. & 0 & 1 & 0.19 & A substantively correct answer. EM scores it zero and F1 scores it 0.19 because the answer is embedded in a full sentence with markup. \\
    \bottomrule
  \end{tabularx}
\end{table}

For a second real example with a stable identifier, MuSiQue item
\path{4hop3__193820_466199_128875_72134} has gold string ``1839'' and the stored
prediction ``The town WLUJ is licensed in, Randolph County, North Carolina, became the capitol of the state where Brantley York was born in 1839.'' Normalisation gives $n(g)=\text{``1839''}$ and the following value for $n(p)$:
\begin{quote}
  \small\ttfamily\raggedright
  town wluj is licensed in randolph county north carolina became capitol of state where brantley york was born in 1839
\end{quote}
Consequently containment is one while EM is zero. This is a provenance example of the implemented lexical rule, not a claim that every statement surrounding the answer is factually correct \citep{trivedi2022musique}.

The choice among these rules is consequential, and the dissertation follows the reproduction protocol of \citet{tran2026single} in reporting \emph{normalised containment} as the headline benchmark figure. The fourth row of \Cref{tab:scoring-examples} shows why. Several of the architectures evaluated here, Debate most of all, answer in a full sentence rather than an answer span, so strict EM records almost none of their correct answers: across the frozen MuSiQue runs at 1,000 tokens, Debate scores 0.030 under EM and 0.510 under containment. An EM-only benchmark would therefore measure output formatting at least as much as reasoning, and would not be comparable across architectures.

Containment buys that comparability at a known cost. It can over-credit an output that mentions an accepted answer and then rejects it, as the third row of \Cref{tab:scoring-examples} shows. It can also under-credit a valid semantic equivalent that is absent from the benchmark's finite reference set. The headline values in \Cref{chap:benchmark-results} are therefore lexical containment rates, not calibrated upper or lower bounds on semantic correctness. \Cref{tab:metric-sensitivity} reports EM and F1 alongside containment so that metric sensitivity is visible. This limitation is also a direct motivation for the rest of the evaluation: a metric that cannot distinguish an endorsed answer from a rejected mention cannot locate a coordination failure. The pipeline additionally records an LLM final-answer judgement using the fixed rubric inherited from the reproduction protocol; it is semantically flexible but model-dependent and can abstain or misparse, and is reported only as a continuity measure. Human labels are reserved for the targeted process claims where the distinction between informational presence and commitment is decisive.

These lexical metrics do not reliably locate intermediate correctness. A worker may express the answer in a sentence rather than an answer field; an agent may use a valid semantic equivalent absent from the alias list; and debate outputs may mention both correct and incorrect candidates. Lexical scores are consequently interpreted as final-output evidence, not as evidence of when the system first reached a correct state.

\section{Semantic stage screen}

Human annotation cannot be applied to the whole trace population, which is why a machine screen exists at all. The five multi-agent architectures produce 26,873 intermediate stage records across 8,000 item-configurations, and each record has to be read in full and labelled on two axes against the question and the accepted answers; that is far beyond what a single annotator can complete. Uniform random sampling is not a workable substitute either, because the events this dissertation is about are rare. Harmful and beneficial final-boundary transitions together account for a low single-digit percentage of items, so a random sample small enough to annotate by hand would contain too few of them to estimate their rates with any precision.

The screen resolves both problems at once. It supplies a complete, if imperfect, labelling of every stage, which is what makes the breadth analysis in \Cref{chap:machine-results} possible; and it partitions the population into transition strata from which rare cases can be deliberately oversampled for human review and then reweighted back to population rates by \Cref{eq:weighting}. The screen is therefore a sampling instrument as much as a measurement in its own right. Because the audit's population estimates depend on the strata it defines, its reliability is tested rather than assumed (\Cref{sec:calibration} and \Cref{chap:human-results}).

A Mistral instruction model was prompted to label intermediate and final stage texts using the reference answer, aliases, question, and numbered stage statements. The frozen judge was Mistral-Small-3.2-24B-Instruct-2506, prompt version \path{stage_discovery_and_endorsement_v4_indexed_evidence}, with temperature 0 and JSON-object output. The final prompt revision applied the same semantic-discovery rule across stages and required evidence-statement identifiers for every correct Discovery label. For each stage, the judge returned both semantic labels, grounded evidence identifiers, an extracted endorsed answer, and reason codes. Outputs that did not parse or lacked a resolvable basis were marked uncertain or parse error rather than silently coerced to correct or wrong. The complete prompt is reproduced in \Cref{app:screening-materials}, together with a full worked example showing the numbered statements a judge receives for one reasoning trace and the JSON it returns.

The full screen covers all 8,000 MAS item-configurations and 26,873 intermediate stages. When the final output text was identical to a recorded stage, its label was reused to avoid unnecessary stochastic re-judging; 2,297 final labels were reused and 5,703 were judged separately. Six uncertain intermediate labels were excluded pairwise from affected transition calculations.

The screen produces four families of descriptive outputs:

\begin{enumerate}[leftmargin=*]
  \item final semantic accuracy, discovery rate, candidate-oracle accuracy, preservation, and aggregation regret for each dataset, model, architecture, and budget;
  \item a mutually exclusive rule-based attribution of terminal failures to no detected Discovery, revision, worker-chain, picker, or aggregator classes;
  \item transitions from Debate proposals to revisions and from revisions to the picker, together with Sequential worker and aggregator transitions;
  \item associations between wall time, realised thinking tokens, truncation, discovery, and post-discovery loss.
\end{enumerate}

The machine screen is explicitly not treated as error-free annotation. It is used to summarise all traces and to create informative sampling strata for human review.

\section{Calibration and the need for a targeted audit}
\label{sec:calibration}

The initial calibration concerned Discovery on 200 stratified stage records, not ordinary final-answer accuracy and not Endorsement. One machine-uncertain record and two Discovery-positive records whose cited evidence could not be located in the evaluated stage were excluded by the evidence-valid rule. On the remaining 197 records, 181 labels agreed with the human labels: 91.9\% accuracy, 98.5\% coverage relative to all 200 records, and Cohen's $\kappa=0.829$. Without the evidence-valid exclusions, the judge decided 199 records and achieved 91.0\% accuracy. Here \emph{grounded} means only that the judge's cited numbered statement can be located in the stage text; it does not independently verify the factual relation, the sufficiency of the oracle passages, or Endorsement correctness. No numerical pass threshold had been preregistered for this earlier calibration, so full-trace use was an operational decision based on the disclosed diagnostics. Performance was not uniform: the stratified sample itself included transition-enriched categories, and Discovery accuracy was 70.0\% for Debate beneficial-flip records and 80.0\% for Sequential \ctow{} records. A targeted sensitivity audit of difficult transitions later found only 73.9\% exact agreement across 199 resolved stage labels ($\kappa=0.478$) and 50.5\% exact agreement across 99 resolved transition pairs (multiclass $\kappa=0.364$).

This difference is not paradoxical. General calibration is dominated by ordinary correct and wrong outputs. Transition-enriched cases place the decision boundary precisely where a judge is most likely to fail: partial evidence, answer mentions without commitment, disagreement between proposals, and truncated conclusions. Because a transition is a compound of two labels, one stage error can change the transition class. The targeted result therefore supersedes the broad calibration result for transition-level inference.

A Qwen2.5-72B-Instruct judge was also evaluated as a stronger alternative. It reached 87.0\% stage accuracy ($\kappa=0.723$) and 79.0\% exact transition accuracy (multiclass $\kappa=0.701$) on the targeted set, with no parse or grounding failures. The later targeted acceptance rule, recorded in \path{analysis/semantic_transition_sensitivity/qwen72b_gptq_v2.score.json}, required stage accuracy $\geq0.90$, stage $\kappa\geq0.75$, transition accuracy $\geq0.80$, transition $\kappa\geq0.75$, and accuracy $\geq0.70$ in every stratum. Qwen failed the first four aggregate thresholds and is not used to relabel the full trace population. This rule is distinct from the earlier operational Discovery calibration. The decision to retain human labels was made from the disclosed reliability criteria rather than from whether machine estimates supported the desired hypothesis.

\section{Human audit scope and sampling}
\label{sec:audit-sampling}

The definitive audit focuses on Debate and Sequential because their final coordination boundaries are well defined and because the machine screen indicated substantively different transition patterns. The frozen machine-labelled population was divided into architecture and transition strata. Rare candidate transitions were oversampled so the audit would include enough difficult cases to diagnose judge direction and component behaviour. Within-stratum selection was fixed before human labels were unblinded.

The final sample contains 160 distinct audited item-configurations drawn from 160 distinct benchmark question identifiers: 80 Debate runs and 80 Sequential runs. Each Debate item-configuration contributes two revised candidates and one final picker output; each Sequential item-configuration contributes the last worker and final aggregator output. This gives 400 non-repeat stage records. Twenty stage records were duplicated under new blind identifiers and inserted as hidden repeats, giving 420 annotation rows in total. The repeat positions and item identities were not visible during annotation.

The audit is a single-annotator study. Hidden repeats therefore measure intra-rater repeatability, not agreement between independent annotators. The sampling design estimates transition rates for the population defined by the frozen machine-screening and sampling snapshot only after stratum weighting; the raw 80-item architecture counts are diagnostic and are not direct prevalence estimates.

\section{Item-level audit schema and the ``15 fields''}

During supervision, the item-level audit columns were described informally as ``15 labels.'' The precise description is \emph{15 item-level audit fields}. They are not fifteen independently annotated error indicators and are not all binary. Four fields are Boolean truncation descriptors, eight are semantic state labels, and three are derived categorical transitions. Consequently, the fields should not be given fifteen separate prevalence interpretations.

At the architecture-aligned source boundary, the Debate source state pools the two revised candidates using an any-correct candidate rule, whereas the Sequential source state is the last worker output. The destination state is the final Debate picker or Sequential aggregator output. \Cref{tab:audit-fields} records the role and scale of each field group.

\begin{table}[htbp]
  \centering
  \small
  \caption[Interpretation of the 15 audit fields]{Interpretation of the 15 named item-level audit fields. The fields are grouped by measurement role rather than treated as fifteen independent outcomes.}
  \label{tab:audit-fields}
  \begin{tabularx}{\textwidth}{@{}>{\raggedright\arraybackslash}p{0.17\textwidth}>{\raggedright\arraybackslash}X>{\raggedright\arraybackslash}p{0.25\textwidth}@{}}
    \toprule
    Group & Fields & Scale and analytical role \\
    \midrule
    Truncation descriptors (4) & \path{question_any_truncated}, \path{source_truncated}, \path{final_truncated}, \path{any_truncated} & Boolean/NA flags for any truncated call in the item-configuration trace, the audited source boundary, the final destination, and their three-valued union. The final-stage metadata are unavailable in this audit snapshot; the union is derived rather than an independent error. \\
    \addlinespace[0.4em]
    Human semantic states (4) & \path{human_source_discovery}, \path{human_source_endorsement}, \path{human_final_discovery}, \path{human_final_endorsement} & Human C/W/U judgements at the source and final stages on the two semantic axes. \\
    \addlinespace[0.4em]
    Human derived transitions (2) & \path{human_endorsement_transition_raw}, \path{human_discovery_to_final_endorsement_raw} & Categorical source-to-final paths. The first is the primary within-axis transition; the second is a cross-axis diagnostic of discovery surviving to final endorsement. \\
    \addlinespace[0.4em]
    Machine semantic states (4) & \path{machine_source_discovery}, \path{machine_source_endorsement}, \path{machine_final_discovery}, \path{machine_final_endorsement} & Frozen Mistral screen labels aligned to the human fields; used for sampling and agreement analysis, not as ground truth. \\
    \addlinespace[0.4em]
    Machine derived transition (1) & \path{machine_transition} & Categorical source-Endorsement-to-final-Endorsement transition used to define the frozen sampling strata. \\
    \bottomrule
  \end{tabularx}
\end{table}

The primary population comparison is therefore based on the weighted categories of \path{human_endorsement_transition_raw}. The cross-axis human transition is reported as a secondary diagnostic; truncation fields are covariates or sensitivity descriptors; and the five machine fields support sampling and human--machine validation. This hierarchy prevents derived or machine-generated columns from being counted as additional human outcomes.

\section{Codebook and annotation procedure}

The codebook was iterated through four revisions during pilot review and frozen before the final unblind analysis. Each record showed the question, accepted gold answer and aliases, evidence context, stage role, and stage text, but concealed architecture, run condition, transition stratum, and whether the row was a repeat. The annotator assigned Discovery and Endorsement independently.

Discovery used the following decision rule:

\begin{itemize}[leftmargin=*]
  \item \correct{} if the stage states the accepted answer or supplies an unambiguous answer-equivalent relationship grounded in the provided evidence;
  \item \wrong{} if no sufficient answer-equivalent content appears, if the relation is wrong, or if only an unsupported option reference is present;
  \item \uncertain{} only when the available text prevents a defensible decision, particularly when truncation cuts through the answer-bearing claim.
\end{itemize}

Endorsement used a stricter commitment rule:

\begin{itemize}[leftmargin=*]
  \item \correct{} if the stage's answer-bearing statement selects the accepted answer;
  \item \wrong{} if it selects a different answer, rejects the accepted answer, or ends without endorsing answer-equivalent content;
  \item \uncertain{} if the commitment itself is missing because the answer-bearing sentence is truncated.
\end{itemize}

The distinction prevents a common evaluator error: crediting a stage merely because the gold string appears somewhere. Evidence fields were checked for every \correct{} label in the final audit. Reason codes were harmonised after unblinding for three records, but the C/W/U labels did not change.

\section{Unblinding and quality checks}

Unblinding occurred after all 420 rows had complete labels. The 20 hidden pairs matched exactly on both axes. Across all rows, Discovery counts were 240 \correct{}, 180 \wrong{}, and zero \uncertain{}; Endorsement counts were 225 \correct{}, 192 \wrong{}, and three \uncertain{}. Fourteen rows had Discovery \correct{} but Endorsement \wrong{}, directly confirming that presence and commitment are empirically distinct. No row had Endorsement \correct{} while Discovery was not \correct{}.

Completeness checks confirmed that all 105 rows labelled Discovery \correct{} in the final audited round contained verbatim or answer-equivalent evidence in the corresponding stage output. The unblind process also corrected a spreadsheet wildcard-count display defect affecting static truncation summary cells. The audited labels and derived transition results were unchanged.

\section{Agreement analysis}

Exact agreement and Cohen's kappa are computed separately for Discovery and Endorsement \citep{cohen1960coefficient}. For hidden repeats, the two blind copies of the same stage are compared. For human-machine agreement, the 400 non-repeat human labels are aligned to the frozen Mistral labels. Percentile bootstrap intervals with resampling at the stage-record level describe uncertainty in kappa \citep{efron1994bootstrap}. Because each audited item contributes two or three source/final stage records, this bootstrap does not adjust for within-item dependence and is reported as a descriptive interval. Agreement direction is also retained: human \correct{}/machine \wrong{} indicates an over-strict machine judgement, while human \wrong{}/machine \correct{} indicates an over-broad judgement.

Transition agreement is more demanding than stage agreement and is reported separately in the targeted sensitivity analysis. A pair is exactly correct only when both source and destination produce the same transition class under human and machine labels. This is why an approximately 80\% stage agreement cannot be interpreted as 80\% reliable transition counts.

\section{Transition estimation and uncertainty handling}

For each audited item, the source and final Endorsement states form one of four primary transitions: \correct$\rightarrow$\correct, \ctow, \wtoc, or \wrong$\rightarrow$\wrong. The three \uncertain{} stage labels are handled under a conservative primary rule that maps \uncertain{} to \wrong{}. Two sensitivity rules are also reported: exclude any item containing \uncertain{}, and map \uncertain{} to \correct{}. The substantive Debate-Sequential harmful-rate ordering is required to remain stable under all three rules.

Population estimates use the frozen stratum sizes as described in \Cref{chap:framework}: 1,185 eligible Debate and 1,037 eligible Sequential item-configurations. The primary harmful quantity is the joint population probability $P(E^{\mathrm{src}}=\correct,E^{\mathrm{final}}=\wrong)$ over the corresponding eligible frame, not a rate conditioned only on correct sources and not a proportion of all 1,600 runs per architecture. A weighted contrast subtracts the Sequential joint harmful rate from the Debate rate. The secondary conditional quantity $P(E^{\mathrm{final}}=\wrong\mid E^{\mathrm{src}}=\correct)$ is obtained from the weighted \ctow{} count divided by the weighted \correct$\rightarrow$\correct plus \ctow{} counts. Approximate 95\% confidence intervals for the primary joint rates use the finite-population stratified variance. Because sample selection included greedy balancing rather than pure random sampling inside every stratum, intervals are treated as descriptive. No claim of conventional randomised-treatment causality is made.

\section{Budget and resource statistics}

Budget outcomes are compared within matched item-configurations using two-sided exact McNemar tests. Discordant counts identify runs that change from wrong to correct and from correct to wrong when the cap rises from 1,000 to 2,000. The pooled tests combine architecture and model conditions, so the same benchmark question can contribute more than one paired observation; their $p$-values do not adjust for this question-level clustering and are interpreted as exploratory. Architecture-level tests reduce but do not remove the same dependence across model variants. Mean paired changes in wall time and realised thinking tokens are reported alongside accuracy differences.

Resource associations use fixed-effect stratification over dataset, model variant, budget, and system. Exposure is defined relative to the within-stratum median for wall time or thinking-token use. Adjusted odds ratios compare semantic discovery and post-discovery drop between high- and low-exposure observations. A separate model tests whether truncation at a destination stage is associated with a \ctow{} transition conditional on a correct source. These analyses are diagnostic. Slower or longer outputs may be consequences of difficult items or unstable trajectories rather than causes of error.

\section{Ethical and data-management considerations}

The experiments use public QA benchmark content and model-generated text; they do not involve recruitment of human participants or collection of personal data. The human annotation is the researcher's analysis of model outputs. The study nevertheless retains an auditable data trail: codebook revisions, blind sheets, round-level completions, the frozen key, and the unblind report are stored separately from the main benchmark outputs. Model-generated statements are not treated as factual sources beyond the supplied benchmark evidence. The analysis reports uncertainty and failed judge acceptance tests rather than concealing them, because evaluator limitations materially affect the validity of the conclusions.
